% !TeX program = xelatex
% !TeX encoding = UTF-8
% Official Linear Algebra final examination booklet.
% Layout adapted from finalexam/main.tex.

\ifdefined\MakeupAnswerVersion
  \documentclass[12pt,a4paper,answers,addpoints]{exam}
\else
  \documentclass[12pt,a4paper,addpoints]{exam}
\fi

\def\CourseTitle{Linear Algebra}
\def\ExamName{Final Examination}
\def\CourseNumber{MATH100}
\def\ExamDate{January, 11th, 2027}
\def\Duration{2}
\def\ExamTime{10:00-12:00}
\def\Term{2609}
\def\ExamDuration{2 Hours}

\usepackage[left=2cm, right=2cm, top=2cm, bottom=2cm]{geometry}

\usepackage{amsmath}
\usepackage{amsthm}
\usepackage{amsfonts}
\usepackage{amssymb}
\usepackage{mathrsfs}
\usepackage{tasks}
% \usepackage[margin=1in]{geometry}
\checkboxchar{$\Box$}
\usepackage{setspace}
%\usepackage{bbding}
%\usepackage{pifont}
\usepackage{wasysym}
\usepackage[shortlabels]{enumitem}
\renewcommand{\qedsymbol}{$\blacksquare$}
\usepackage{array} 
\usepackage{multicol}
\setlength{\columnsep}{1cm}
\usepackage{afterpage}
\usepackage{color} 
\usepackage{graphicx} % in preamble
\usepackage{listings}
\usepackage{upquote}

\lstset{
    escapeinside={(*@}{@*)},           % Allow LaTeX commands between (*@ and @*)
}

% new macros
\newcommand{\todo}[1]{\textcolor{red}{(To do): #1}}
\newcommand{\tored}[1]{\textcolor{red}{\textbf{#1}}}
\usepackage{xeCJK}
\defaultfontfeatures{Ligatures=TeX}

\setCJKmainfont[
    BoldFont={Noto Sans TC}
]{Noto Sans TC}

% Required by listings/typewriter text
\setCJKmonofont{Noto Sans TC}

\newcommand\blankpage{%
    \null
    \thispagestyle{empty}%
    \addtocounter{page}{-1}%
    \shownoanswers{\newpage}}
%% https://tex.stackexchange.com/a/203986
% Create a True-False question format
\newcommand*{\TrueFalse}[1]{%
\ifprintanswers
    \ifthenelse{\equal{#1}{T}}{%
        \textbf{TRUE}\hspace*{14pt}False
    }{
        True\hspace*{14pt}\textbf{FALSE}
    }
\else
    {True}\hspace*{20pt}False
\fi
} 
%% The following code is based on an answer by Gonzalo Medina
%% https://tex.stackexchange.com/a/13106/39194
\newlength\TFlengthA
\newlength\TFlengthB
\settowidth\TFlengthA{\hspace*{1.16in}}
\newcommand\TFQuestion[2]{%
    \setlength\TFlengthB{\linewidth}
    \addtolength\TFlengthB{-\TFlengthA}
    \parbox[t]{\TFlengthA}{\TrueFalse{#1}}\parbox[t]{\TFlengthB}{#2}}
%%

%%
\pagestyle{headandfoot}
%\extraheadheight{.25in}
\runningheader{{\CourseNumber} - {\CourseTitle}}
{}
{{\ExamName} - {\Term}} 
\firstpagefooter{}{Page \thepage\ of \numpages}{Please go on to the next page\ldots} %
\runningfooter{}{Page \thepage\ of \numpages}{\iflastpage{End of exam.}{Please go on to the next page\ldots}}

\renewcommand{\fillinlinelength}{1cm}  % change default length

\renewenvironment{solution}
  {\par\noindent\textbf{Solution: }\color{red}} % begin code
  {\par}    

% --- show solution chunks only when 'answers' is active ---
\usepackage{xcolor} % for \color

% Inline macro: use as \showsol{<solution text>}
\newcommand{\showsol}[1]{\ifprintanswers{\textbf{\color{red}\underline{Solution:} }\color{red}#1}\fi}

\newenvironment{show_code_sol}
  {\par\noindent\textbf{Solution:}\par
   \lstset{basicstyle=\ttfamily\small\color{red}}}
  {\par}


\usepackage{listings}
\usepackage{comment} % for \excludecomment (safe with verbatim)
% \usepackage[T1]{fontenc}


% Environment version that mimics 'solution' but guarded
\newenvironment{solonly}
  {\ifprintanswers\par\noindent\textbf{Solution: }\color{red}}
  {\par\fi}

% If you want to keep using the built-in 'solution' name,
% make sure it stays conditional:
\renewenvironment{solution}
  {\ifprintanswers\par\noindent\textbf{Solution: }\color{red}}
  {\par\fi}

\newcommand{\shownoanswers}[1]{\ifprintanswers\else #1\fi}

% Base style for C
% One listings + comment load
\usepackage{listings}
\usepackage{comment} % safe with verbatim

% --- styles (no frame, no line numbers) ---
\lstdefinestyle{examc}{
  language=C,
  basicstyle=\ttfamily\small,
  numbers=none,
  frame=none,
  columns=fullflexible,
  keepspaces=true,
  showstringspaces=false,
  tabsize=4,
  breaklines=true,
  keywordstyle=\bfseries,
  commentstyle=\itshape,
  stringstyle=\ttfamily
}
\lstdefinestyle{examc-sol}{
  style=examc,
  basicstyle=\ttfamily\small\color{red}
}

% --- Environments ---
\lstnewenvironment{Ccode}[1][]
  {\lstset{style=examc,#1}}
  {}

% Define CcodeSol only when answers are printed; otherwise exclude it.
\ifprintanswers
  \lstnewenvironment{CcodeSol}[1][]
    {\lstset{style=examc-sol,#1}}
    {}
\else
  \excludecomment{CcodeSol} % body is ignored; no verbatim processing
\fi

% Inline C
\newcommand{\Cinline}[1]{\lstinline[language=C]{#1}}

% ---------------------------------------------------------------------
  


\begin{document}

% \noindent
% \begin{minipage}[t]{0.48\textwidth}
%     \vspace{-10pt}
%     \fbox{\parbox[c][1.2cm][c]{4cm}{Seat\\No.}}
% \end{minipage}

% \vspace{-2cm}
% \begin{center}
%     \includegraphics[width=0.27\textwidth]{must-logo.pdf}
% \end{center}

% \begin{center}
%     \vspace*{0.5cm}
%     {\textbf{\Huge 澳\hspace{2em}門\hspace{2em}科\hspace{2em}技\hspace{2em}大\hspace{2em}學}}\\[0.3cm]
%     {\large \textbf{MACAU UNIVERSITY OF SCIENCE AND TECHNOLOGY}}\\[1cm]

%     {\Large \textbf{School of Computer Science and Engineering}}\\[0.2cm]
%     {\Large \textbf{Faculty of Innovation Engineering}}\\[0.2cm]
%     {\Large \textbf{Bachelor of Science}}\\[1cm]

%     {\large \textbf{2609}}\\[1cm]
%     {\Large \textbf{\CourseNumber \quad \CourseTitle}}\\[1cm]
% \end{center}

% \begin{center}
% \begin{tabular}{@{}p{4cm}p{8cm}@{}}
%     \textbf{No. of Pages:} & {\numpages} Pages (including this page)\\[0.5cm]
%     \textbf{Exam Date:} & \ExamDate\\[0.5cm]
%     \textbf{Exam Duration:} & \ExamDuration\ (\ExamTime)\\[0.5cm]
%     \textbf{Student Name:} & \underline{\hspace{7cm}}\\[0.5cm]
%     \textbf{Student Number:} & \underline{\hspace{7cm}}\\
% \end{tabular}
% \end{center}

% \vspace{0.5cm}

% \noindent
% \textbf{Please read the following instructions carefully.}
% \begin{itemize}[leftmargin=*]
%     \item This question booklet contains \textbf{\numquestions} questions and
%           \textbf{\numpages} pages (including the cover), for a total of
%           \textbf{\numpoints} marks. Check that all pages are present.
%     \item Return this question booklet, together with all extra unused answer sheets
%           and draft paper, at the end of the examination.
%     \item This is a \textbf{Closed Book Examination}. Dictionaries, calculators,
%           mobile phones, smart watches, and other electronic devices are
%           \textbf{NOT permitted}.
%     \item Show sufficient work in Part III. Unless instructed otherwise, give exact answers.
%     \item \textbf{Answers must be provided in English to receive marks.}
% \end{itemize}

% \vfill

\begin{questions}
\addpoints

\newpage
\fullwidth{\textbf{\Large Part I. Fill in the blanks.}
\hfill\textbf{(20 marks; 2 marks each)}}
\vspace{0.4cm}

\question[2]
Consider the basis
$\left\{\begin{bmatrix}1\\2\end{bmatrix},
\begin{bmatrix}2\\-1\end{bmatrix}\right\}$ for $\mathbb R^2$.
The coordinate vector of $\begin{bmatrix}5\\0\end{bmatrix}$ with respect to this basis is
$\underline{\hspace{28mm}}$.
\showsol{$\begin{bmatrix}1\\2\end{bmatrix}$.}
\showoutcome{3}{a (a1)}

\question[2]
If $A$ is invertible and $\det(A)=\alpha$, then $\det(A^{-1})=\underline{\hspace{35mm}}$.
\showsol{$\dfrac{1}{\alpha}$.}
\showoutcome{2}{a (a1)}

\question[2]
If $A=\begin{bmatrix}2&-1\\4&3\end{bmatrix}$, then
$A^{-1}=\underline{\hspace{35mm}}$.
\showsol{$A^{-1}=\begin{bmatrix}0.3&0.1\\-0.4&0.2\end{bmatrix}$.}
\showoutcome{2}{a (a1)}

\question[2]
If $A$ has 7 columns and $\operatorname{rank}(A)=4$, then the dimension of the
null space of $A$ is $\underline{\hspace{22mm}}$.
\showsol{$3$.}
\showoutcome{3}{a (a1)}

\question[2]
For an eigenvalue $\lambda$, the corresponding eigenspace is the null space of
\par\hspace*{8mm}$\underline{\hspace{30mm}}$.
\showsol{$A-\lambda I$.}
\showoutcome{6}{a (a2)}

\question[2]
The eigenvalues of $\begin{bmatrix}3&1\\0&-2\end{bmatrix}$ are
$\underline{\hspace{35mm}}$.
\showsol{$3$ and $-2$.}
\showoutcome{6}{a (a2)}

\question[2]
$\det\!\begin{bmatrix}4&-2\\1&3\end{bmatrix}=\underline{\hspace{24mm}}$.
\showsol{$14$.}
\showoutcome{2}{a (a1)}

\question[2]
$\begin{bmatrix}-2&3&6\end{bmatrix}\times\begin{bmatrix}-4\\3\\6\end{bmatrix}
=\underline{\hspace{24mm}}$.
\showsol{$7$.}
\showoutcome{53}{a (a1)}

\question[2]
The distance between $\begin{bmatrix}1\\-2\\3\end{bmatrix}$ and
$\begin{bmatrix}-2\\2\\3\end{bmatrix}$ is $\underline{\hspace{24mm}}$.
\showsol{$5$.}
\showoutcome{7}{a (a1)}

\question[2]
If $A=\begin{bmatrix}2&0&0\\0&-3&0\\0&0&4\end{bmatrix}$, then
$A^2=\underline{\hspace{42mm}}$.
\showsol{$A^2=\operatorname{diag}(4,9,16)$.}
\showoutcome{5}{a (a2)}

\newpage
\fullwidth{\textbf{\large Part II. Select the only correct statement.}
\hfill\textbf{(20 marks; 2 marks each)}}
\vspace{0.3cm}

\question[2]
For a square matrix $A$, which condition is equivalent to $A$ being invertible?
\begin{choices}
    \choice $A\mathbf x=\mathbf 0$ has a nontrivial solution.
    \choice $A\mathbf x=\mathbf b$ has a unique solution for every $\mathbf b$.
    \choice The columns of $A$ are linearly dependent.
    \choice $\det(A)=0$.
\end{choices}
\showcorrectanswer{B. $A\mathbf x=\mathbf b$ has a unique solution for every $\mathbf b$.}
\showoutcome{1}{a (a2)}

\question[2]
If $B$ is obtained from $A$ by interchanging two rows, then
\begin{choices}
    \choice $\det(B)=\det(A)$.
    \choice $\det(B)=-\det(A)$.
    \choice $\det(B)=2\det(A)$.
    \choice $\det(B)=0$ for every $A$.
\end{choices}
\showcorrectanswer{B. $\det(B)=-\det(A)$.}
\showoutcome{2}{a (a1)}

\question[2]
Which set is a subspace of $\mathbb R^3$?
\begin{choices}
    \choice $\{(x,y,z):x-2y+z=0\}$
    \choice $\{(x,y,z):x-2y+z=1\}$
    \choice $\{(x,y,z):x,y,z>0\}$
    \choice $\{(x,y,z):x+y+z\geq 0\}$
\end{choices}
\showcorrectanswer{A. $\{(x,y,z):x-2y+z=0\}$.}
\showoutcome{4}{a (a1)}

\question[2]
Pivot columns 1 and 3 occur in the reduced echelon form of $A$. A basis for
$\operatorname{Col}(A)$ is obtained from
\begin{choices}
    \choice columns 1 and 3 of the original matrix $A$.
    \choice columns 1 and 3 of the reduced echelon form of $A$.
    \choice rows 1 and 3 of $A$.
    \choice all nonzero rows of the reduced echelon form of $A$.
\end{choices}
\showcorrectanswer{A. Columns 1 and 3 of the original matrix $A$.}
\showoutcome{3}{a (a1)}

\question[2]
If $A$ has 5 columns and $\operatorname{rank}(A)=3$, then the dimension of the
null space of $A$ equals
\begin{choices}
    \choice 1
    \choice 2
    \choice 3
    \choice 5
\end{choices}
\showcorrectanswer{B. $2$.}
\showoutcome{3}{a (a1)}

\newpage

\begingroup
\small

\question[2]
Which statement about the fundamental subspaces of an $m\times n$ matrix $A$ is correct?
\begin{choices}
    \choice $\operatorname{Null}(A)$ is a subspace of $\mathbb R^m$.
    \choice The pivot columns of the original matrix $A$ form a basis for
            $\operatorname{Col}(A)$.
    \choice Row-equivalent matrices always have the same column space.
    \choice The system $A\mathbf x=\mathbf b$ is inconsistent exactly when
            $\mathbf b\in\operatorname{Col}(A)$.
\end{choices}
\showcorrectanswer{B. The pivot columns of the original matrix $A$ form a basis for
$\operatorname{Col}(A)$.}
\showoutcome{3}{a (a1)}

\vspace{2mm}

\question[2]
If $A$ is a $4\times7$ matrix, then the rank of $A^T$ is at most
\begin{choices}
    \choice 7
    \choice 3
    \choice 4
    \choice 8
\end{choices}
\showcorrectanswer{C. $4$.}
\showoutcome{2}{a (a1)}

\vspace{2mm}

\question[2]
If $A$, $B$, and $C$ are square matrices of the same order, then $(ABC)^T$ equals
\begin{choices}
    \choice $C^TB^TA^T$
    \choice $A^TB^TC^T$
    \choice $C^TA^TB^T$
    \choice $(CBA)^T$
\end{choices}
\showcorrectanswer{A. $C^TB^TA^T$.}
\showoutcome{5}{a (a2)}

\vspace{2mm}

\question[2]
Let $A$ be a diagonalizable $5\times5$ matrix. Which statement must be true?
\begin{choices}
    \choice There is a basis of $\mathbb R^5$ consisting of eigenvectors of $A$.
    \choice The matrix $A$ has five distinct eigenvalues.
    \choice The matrix $A$ must be symmetric.
    \choice Any five eigenvectors of $A$ are linearly dependent.
\end{choices}
\showcorrectanswer{A. There is a basis of $\mathbb R^5$ consisting of eigenvectors of $A$.}
\showoutcome{6}{a (a2)}

\vspace{2mm}

\question[2]
Let $A$ be an $n\times n$ matrix. The equation $A\mathbf x=6\mathbf x$ has only
the trivial solution if and only if which matrix is invertible?
\begin{choices}
    \choice $A$
    \choice $A+6I$
    \choice $A-6I$
    \choice $6A$
\end{choices}
\showcorrectanswer{C. $A-6I$.}
\showoutcome{6}{a (a2)}

\endgroup

\vspace{0.3cm}
\fullwidth{\textbf{End of Part II. Continue to Part III and show all work.}}

\newpage
\fullwidth{\textbf{\large Part III. Show all your work in the space provided.}
\hfill\textbf{(60 marks)}}
\vspace{0.4cm}

\question[12]
Use Gaussian elimination to solve the system
\[
\begin{aligned}
x+y+z&=2,\\
2x-y+z&=-1,\\
x+2y-z&=6.
\end{aligned}
\]
State whether the solution is unique and justify your conclusion.

\shownoanswers{\vspace{8cm}}
\showsol{Using elementary row operations,
\[
\left[\begin{array}{ccc|c}
1&1&1&2\\
2&-1&1&-1\\
1&2&-1&6
\end{array}\right]
\sim
\left[\begin{array}{ccc|c}
1&1&1&2\\
0&1&-2&4\\
0&0&-7&7
\end{array}\right].
\]
Thus $z=-1$, $y=2$, and $x=1$. The solution is unique because there is a
pivot in every variable column.}
\showoutcome{1}{a (a2)}

% \newpage

\thispagestyle{headandfoot}
\question
Let
\[
A=\begin{bmatrix}
2&-1&3\\
1&2&4\\
3&1&7
\end{bmatrix}.
\]
\begin{parts}
    \part[4] Find a row echelon form of $A$ and identify the pivot columns.\par\vspace{2mm}
    \part[4] Find a basis for $\operatorname{Null}(A)$.\par\vspace{2mm}
    \part[4] Find a basis for $\operatorname{Col}(A)$.\par\vspace{2mm}
    % \part[3] Verify the rank-nullity theorem for $A$.
\end{parts}

\shownoanswers{\vspace{9cm}}
\showsol{\textbf{(a)} One row echelon form is
$\begin{bmatrix}1&2&4\\0&-5&-5\\0&0&0\end{bmatrix}$, so the pivot columns are
columns 1 and 2.}
\showoutcome{2}{a (a1)}

\showsol{\textbf{(b)} Solving $A\mathbf x=\mathbf0$ gives
$\mathbf x=t\begin{bmatrix}-2\\-1\\1\end{bmatrix}$. Hence a basis for
$\operatorname{Null}(A)$ is
$\left\{\begin{bmatrix}-2\\-1\\1\end{bmatrix}\right\}$.}
\showoutcome{3}{a (a1)}

\showsol{\textbf{(c)} A basis for $\operatorname{Col}(A)$ is formed from the
corresponding pivot columns of the original matrix:
$\left\{\begin{bmatrix}2\\1\\3\end{bmatrix},
\begin{bmatrix}-1\\2\\1\end{bmatrix}\right\}$.}
\showoutcome{3}{a (a1)}

\newpage
\thispagestyle{headandfoot}
\question
Consider
\[
B=\begin{bmatrix}
2&1&-1\\
-1&4&1\\
-2&2&3
\end{bmatrix}.
\]
\begin{parts}
    \part[5] Find the eigenvalues of $B$ and a basis for each eigenspace.\par\vspace{2mm}
    \part[4] Diagonalize $B$ by finding matrices $P$ and $D$ such that
             $B=PDP^{-1}$.\par\vspace{2mm}
    \part[3] Use the diagonalization to compute $B^3$.
\end{parts}

\shownoanswers{\vspace{7cm}}
\showsol{\textbf{(a)} The characteristic polynomial is
$\det(\lambda I-B)=(\lambda-1)(\lambda-3)(\lambda-5)$. Bases for the
eigenspaces are
\[
E_1=\operatorname{span}\left\{\begin{bmatrix}1\\0\\1\end{bmatrix}\right\},\quad
E_3=\operatorname{span}\left\{\begin{bmatrix}1\\1\\0\end{bmatrix}\right\},\quad
E_5=\operatorname{span}\left\{\begin{bmatrix}0\\1\\1\end{bmatrix}\right\}.
\]}
\showoutcome{6}{a (a2)}

\showsol{\textbf{(b)} Taking the eigenvectors as columns gives
\[
P=\begin{bmatrix}1&1&0\\0&1&1\\1&0&1\end{bmatrix},\qquad
D=\begin{bmatrix}1&0&0\\0&3&0\\0&0&5\end{bmatrix},\qquad
P^{-1}=\frac12\begin{bmatrix}1&-1&1\\1&1&-1\\-1&1&1\end{bmatrix}.
\]
Therefore $B=PDP^{-1}$.}
\showoutcome{6}{a (a2)}

\showsol{\textbf{(c)} Since $B^3=PD^3P^{-1}$,
\[
B^3=\begin{bmatrix}
14&13&-13\\
-49&76&49\\
-62&62&63
\end{bmatrix}.
\]}
\showoutcome{6}{a (a2)}

\newpage
\thispagestyle{headandfoot}
\question[12]
Compute the determinant of the matrix
\[
M=\begin{bmatrix}
0&1&1&1&1\\
2&3&3&3&3\\
2&3&5&5&5\\
2&3&5&7&7\\
2&3&5&7&11
\end{bmatrix}.
\]


\shownoanswers{\vspace{9cm}}
\showsol{Subtract each row from the row below it, working from the bottom upward,
and then replace $R_2$ by $R_2-3R_1$. This gives
\[
\begin{bmatrix}
0&1&1&1&1\\
2&0&0&0&0\\
0&0&2&2&2\\
0&0&0&2&2\\
0&0&0&0&4
\end{bmatrix}.
\]
Interchanging the first two rows produces an upper triangular matrix. Therefore
\[
\det(M)=-(2)(1)(2)(2)(4)=-32.
\]}
\showoutcome{2}{a (a1)}

\newpage
\thispagestyle{headandfoot}
\question
Let $a,b\in\mathbb R$. Consider the system
\[
\begin{aligned}
x+2y+z&=a,\\
2x+4y+3z&=b,\\
3x+6y+4z&=2a-b.
\end{aligned}
\]
\begin{parts}
    \part[6] Determine all values of $a$ and $b$ for which the system is consistent.\par\vspace{2mm}
    \part[6] Are there any values of $a$ and $b$ for which the system has exactly one
             solution? Justify your answer.
\end{parts}

\shownoanswers{\vspace{5cm}}
\showsol{\textbf{(a)} The third coefficient row is the sum of the first two.
Consistency therefore requires the same relation on the right-hand side:
$2a-b=a+b$, or $a=2b$. For these values, with $y=t$, the solutions are
\[
(x,y,z)=(5b-2t,\ t,\ -3b),\qquad t\in\mathbb R.
\]}
\showoutcome{1}{a (a2)}

\showsol{\textbf{(b)} No. Whenever the system is consistent, the coefficient
matrix has rank $2$ but there are $3$ unknowns. Hence one variable is free and
the system has infinitely many solutions, never exactly one solution.}
\showoutcome{1}{a (a2)}

% \newpage
% \thispagestyle{headandfoot}
% \fullwidth{\textbf{Additional workspace for Question 25}}
% \vfill

% \begin{center}
%     \Large\textbf{END OF EXAMINATION}
% \end{center}

\end{questions}

\vfill
\noindent\makebox[2.5in]{\textbf{\Large Score:} \hrulefill}

\end{document}
